\documentclass[
   12pt,
   oneside,
   a4paper,
   english,
   brazil
]{abntex2}

% Codificação e fontes
\usepackage[utf8]{inputenc}
\usepackage[T1]{fontenc}
\usepackage{lmodern}

% Microtipografia
\usepackage{microtype}

% Idioma
\usepackage[english, brazil]{babel}

% Citações e aspas tipográficas
\usepackage{csquotes}

% Figuras e gráficos
\usepackage{graphicx}
\graphicspath{{sections/figs/}}
\usepackage{caption}
\usepackage{subcaption}

% Matemática
\usepackage{amsmath, amssymb, mathtools}

% Tabelas melhoradas
\usepackage{booktabs}
\usepackage{array}

% Listas
\usepackage{enumitem}
\setlist{nosep}

% Espaçamento inicial
\usepackage{indentfirst}

% Citações ABNT
\usepackage{abntex2cite}
\bibliographystyle{abntex2-alf}

% Hiperlinks
\usepackage{hyperref}
\hypersetup{
  pdftitle={AUTOMATIC DETECTION OF EQUATORIAL PLASMA BUBBLES IN AIRGLOW IMAGES USING MACHINE LEARNING TECHNIQUES},
  pdfauthor={Alberto Ivo da Costa Vieira},
  pdfsubject={Qualificação de Mestrado},
  pdfkeywords={Equatorial Plasma Bubbles, Artificial Intelligence, Machine Learning, Deep Learning, Space Weather, EMBRACE, INPE},
  colorlinks=true,
  linkcolor=black,
  citecolor=black,
  urlcolor=blue
}

% Metadados / Capa / Folha de rosto
\titulo{AUTOMATIC DETECTION OF EQUATORIAL PLASMA BUBBLES IN AIRGLOW IMAGES USING MACHINE LEARNING TECHNIQUES}
\autor{Alberto Ivo da Costa Vieira}
\local{João Pessoa-PB}
\data{2026}
\instituicao{
  Universidade Federal da Paraíba
  \par
  Programa de Pós-Graduação em Informática
}
\tipotrabalho{Master's Qualification Exam}
\preambulo{Qualification Exam presented to the Graduate Program in Informatics of the Universidade Federal da Paraíba as a partial requirement for obtaining the degree of Master in Informatics.}

\begin{document}

\selectlanguage{english}

% ================= PRE-TEXTUAL =================
\pretextual

\imprimircapa
\imprimirfolhaderosto*

% --- Dedicatória ---
\begin{dedicatoria}
   \vspace*{\fill}
   \centering
   \noindent
   \textit{This work is dedicated to my family, \\
   friends, and everyone who supported me \\
   throughout this journey.} \vspace*{\fill}
\end{dedicatoria}

% --- Agradecimentos ---
\begin{agradecimentos}
   \noindent
   I would like to express my sincere gratitude to my advisor and co-advisor for their continuous support and guidance.   

   \noindent I also thank my colleagues and the researchers at INPE for providing the data and expertise necessary for this work.
\end{agradecimentos}

% --- Epígrafe ---
\begin{epigrafe}
    \vspace*{\fill}
    \begin{flushright}
        \textit{``Science is a way of thinking much more than it is a body of knowledge.''\\
        (Carl Sagan)}
    \end{flushright}
\end{epigrafe}

% --- Resumo / Abstract ---
% The abstract is inputted from sections/abstract.tex
\begin{resumo}[Abstract]
\begin{otherlanguage*}{english}
Equatorial Plasma Bubbles (EPBs) trigger severe ionospheric scintillations, often disrupting satellite communication and navigation. For years, researchers have relied on manual detection in all-sky airglow imagery. Though foundational, this approach requires immense labor and remains sensitive to environmental noise, pointing to a clear need for automation. Deep learning has certainly opened new doors here, but throwing raw optical imagery directly into complex networks tends to spike computational costs and amplify background artifacts. In response, we developed a physics-informed machine learning architecture tailored specifically for EPB detection. The pipeline begins with domain-specific preprocessing. By modeling and subtracting diffuse background illumination, alongside Contrast Limited Adaptive Histogram Equalization (CLAHE), we naturally accentuate the physical plasma structures. Rather than flattening these images and losing their spatial context, we applied Two-Dimensional Principal Component Analysis (2DPCA) to extract features directly. This step substantially reduces dimensionality while keeping the two-dimensional spatial correlations intact. A stacking ensemble classifier, combining Support Vector, XGBoost, and Random Forest algorithms, evaluates these features. We also incorporated SMOTE to manage the inherent class imbalance of clear skies versus bubble events. Across a 5-fold cross-validation, the model reached an average Accuracy and F1-Score of 99.80\%. Beyond just metrics, SHapley Additive exPlanations (SHAP) proved that the algorithm actually focuses on the physical morphology of the depletions, ignoring random environmental noise. When tested on continuous nocturnal sequences, the system maintained strong stability with virtually no false positives, offering a fast and highly interpretable tool for real-time space weather monitoring.

\vspace{\onelineskip}
\noindent \textbf{Keywords:} Equatorial Plasma Bubbles. EPB. Machine Learning. Space Weather. All-Sky Imager. 2DPCA. Stacking Ensemble. Ionospheric Irregularities.
\end{otherlanguage*}
\end{resumo}

\begin{resumo}[Resumo]
\begin{otherlanguage*}{brazil}
Bolhas de Plasma Equatoriais (EPBs) desencadeiam cintilações ionosféricas severas, frequentemente interrompendo a comunicação e a navegação por satélite. Durante anos, os pesquisadores confiaram na detecção manual em imagens de luminescência atmosférica de todo o céu. Embora fundamental, essa abordagem exige um trabalho imenso e permanece sensível ao ruído ambiental, apontando para uma clara necessidade de automação. O aprendizado profundo certamente abriu novas portas aqui, mas o uso de imagens ópticas brutas diretamente em redes complexas tende a aumentar os custos computacionais e amplificar artefatos de fundo. Em resposta, desenvolvemos uma arquitetura de aprendizado de máquina baseada na física, adaptada especificamente para a detecção de EPB. O processo começa com um pré processamento específico do domínio. Ao modelar e subtrair a iluminação de fundo difusa, juntamente com a Equalização de Histograma Adaptativa Limitada por Contraste (CLAHE), nós acentuamos naturalmente as estruturas físicas do plasma. Em vez de achatar essas imagens e perder seu contexto espacial, aplicamos a Análise de Componentes Principais Bidimensional (2DPCA) para extrair características diretamente. Essa etapa reduz substancialmente a dimensionalidade, mantendo intactas as correlações espaciais bidimensionais. Um classificador ensemble stacking, combinando os algoritmos Support Vector, XGBoost e Random Forest, avalia essas características. Também incorporamos o SMOTE para gerenciar o desequilíbrio de classe inerente entre céus limpos e eventos de bolhas. Ao longo de uma validação cruzada de 5 dobras, o modelo atingiu uma Acurácia média e um F1 Score de 99.80\%. Além de apenas métricas, as explicações aditivas de SHapley (SHAP) provaram que o algoritmo realmente foca na morfologia física das depleções, ignorando ruídos ambientais aleatórios. Quando testado em sequências noturnas contínuas, o sistema manteve forte estabilidade virtualmente sem falsos positivos, oferecendo uma ferramenta rápida e altamente interpretável para o monitoramento do clima espacial em tempo real.

\vspace{\onelineskip}
\noindent Palavras chave: Bolhas de Plasma Equatoriais. EPB. Aprendizado de Máquina. Clima Espacial. Imageador de Todo o Céu. 2DPCA. Ensemble Stacking. Irregularidades Ionosféricas.
\end{otherlanguage*}
\end{resumo}


% --- Listas ---
\listoffigures*
\cleardoublepage

\listoftables*
\cleardoublepage

% --- Siglas ---
\begin{siglas}
  \item[2DPCA] Two-Dimensional Principal Component Analysis
  \item[CNN] Convolutional Neural Network
  \item[DL] Deep Learning
  \item[EMBRACE] Estudo e Monitoramento Brasileiro de Clima Espacial
  \item[EPB] Equatorial Plasma Bubble
  \item[GNSS] Global Navigation Satellite Systems
  \item[INPE] Instituto Nacional de Pesquisas Espaciais
  \item[IoU] Intersection over Union
  \item[ML] Machine Learning
  \item[SVM] Support Vector Machine
  \item[XAI] Explainable Artificial Intelligence
\end{siglas}

% --- Sumário ---
\tableofcontents*
\cleardoublepage

% ================= TEXTUAL =================
\textual

% --- Capítulos ---
\chapter{Introduction}\label{chap:introduction}

Equatorial Plasma Bubbles (EPBs) are a major space weather hazard triggered by the Rayleigh-Taylor instability in the post-sunset equatorial ionosphere. Because these large-scale plasma depletions cause significant amplitude and phase scintillations in trans-ionospheric radio signals, they can severely degrade Global Navigation Satellite Systems (GNSS) and satellite communications. With our growing reliance on these technologies, continuous and real-time monitoring of EPBs has never been more important \cite{zhong2025automatic,okoh2025bootstrapping,githio2024machine}.

For decades, tracking these bubbles meant relying heavily on All-Sky Imagers (ASIs) and having experts manually sift through vast amounts of high-resolution visual data \cite{okoh2025bootstrapping,zhong2025automatic}. Though this human-centric approach laid the groundwork for our current understanding, it demands immense labor and remains quite vulnerable to environmental interference, like shifting lunar illumination or passing clouds. Today, the sheer volume of observational data makes manual processing a bottleneck. Machine Learning (ML) stands out as a clear alternative to automate this task. However, feeding raw optical imagery directly into standard networks often drives up computational costs and forces models to battle background noise, highlighting the need for smarter, physics-aware detection strategies \cite{githio2024machine}.

In this work, we propose a machine learning architecture designed specifically around the physics of all-sky airglow imagery. Rather than treating the images simply as generic grids of pixels, the pipeline applies targeted preprocessing to strip away non-uniform illumination and boost local contrast, followed by a dimensionality reduction strategy that preserves the natural two-dimensional structure of the data. An ensemble classifier then evaluates the extracted features, merging the predictive strengths of several base algorithms to identify plasma structures reliably.

Ultimately, our goal is to validate a streamlined detection system for EPBs that respects the underlying physics of the data while remaining interpretable, fast, and computationally light enough for real-time monitoring.

The remainder of this document is organized as follows. Chapter~\ref{chap:objectives} presents the general and specific objectives that guide this research. Chapter~\ref{chap:methodology} describes the full technical methodology, covering dataset acquisition, the physics-informed preprocessing pipeline, feature extraction via 2DPCA, the stacking ensemble classifier, and the evaluation strategy. Chapter~\ref{chap:results} reports the preliminary results obtained so far, including cross-validation performance, model interpretability analysis, and temporal stability tests. Chapter~\ref{chap:schedule} outlines the execution schedule for the remaining phases of the dissertation. Finally, Chapter~\ref{chap:conclusion} summarizes the main findings and discusses directions for future work.
\chapter{Objectives}\label{chap:objectives}

This chapter presents the objectives that guide the development of this dissertation. They are organized into a general objective, which defines the overarching research goal, and a set of specific objectives that outline the individual steps required to achieve it.

\section{General Objective}

The primary goal of this research is to develop and evaluate machine learning models capable of automating the detection of Equatorial Plasma Bubbles in airglow images at the OI 630 nm emission from the EMBRACE INPE program, contributing to the improvement of operational space weather monitoring in Brazil.

\section{Specific Objectives}

To achieve this goal, the following specific objectives guide the development of this dissertation:

\begin{enumerate}
    \item Conduct a systematic literature review on the application of machine learning techniques for detecting ionospheric irregularities and structures in airglow imagery, identifying the architectures, preprocessing strategies, and evaluation metrics currently in use.
    \item Build a labeled database of all-sky airglow images from the EMBRACE INPE archives, containing examples of both EPB presence and absence across different observational conditions, seasons, and levels of solar activity.
    \item Design and implement a physics-informed preprocessing pipeline that models and removes background illumination, enhances local contrast through CLAHE, and normalizes pixel intensities to isolate the physical signatures of plasma depletions.
    \item Extract spatial features from the preprocessed images using Two-Dimensional Principal Component Analysis (2DPCA), preserving the two-dimensional correlations while substantially reducing data dimensionality.
    \item Train and evaluate a stacking ensemble classifier combining Support Vector, XGBoost, and Random Forest algorithms, integrated with SMOTE to address class imbalance, for the automatic classification of EPB events.
    \item Assess the interpretability of the trained model through SHapley Additive exPlanations (SHAP), verifying that the classifier decisions are driven by the physical morphology of plasma depletions rather than environmental artifacts.
    \item Validate the temporal stability of the system over continuous nocturnal observation sequences, ensuring reliable performance under real-time monitoring conditions.
    \item Compare the proposed architecture against alternative machine learning approaches using quantitative metrics such as precision, recall, F1 score, and Intersection over Union, complemented by domain expert validation.
\end{enumerate}
\chapter{Methodology}\label{chap:methodology}

The methodology for the automated detection of Equatorial Plasma Bubbles (EPBs) integrates physics-informed image preprocessing, advanced feature extraction, and ensemble machine learning. This pipeline is built to reliably spot EPB signatures in all-sky airglow imagery while naturally filtering out environmental noise and class imbalance.

\subsection{Dataset Acquisition}
All optical imagery used in this research was sourced directly from the EMBRACE (Estudo e Monitoramento Brasileiro do Clima Espacial) project, which is maintained by the Brazilian National Institute for Space Research (INPE - Instituto Nacional de Pesquisas Espaciais). This dataset gives us a view of ionospheric conditions across the Brazilian sector, capturing everything from perfectly clear nights to intense EPB events.

\subsection{Physics-Informed Image Preprocessing}
Raw optical images of the night sky rarely offer perfect viewing conditions; they are often skewed by shifting lunar phases or sensor gradients. To counter this, we apply a domain-specific preprocessing pipeline before any machine learning takes place. First, we model and remove the background illumination using a large Gaussian filter ($\sigma=10$). Subtracting this diffuse modeled light from the original image effectively clears out environmental biases. Next, we use Contrast Limited Adaptive Histogram Equalization (CLAHE) with a clip limit of 0.03. This localized tweak makes the dark, vertical tears of plasma depletions stand out sharply against the night sky. Finally, normalizing all pixel intensities to a standard $[0, 1]$ range keeps the math stable during feature extraction.

\subsection{Data Augmentation and Balancing}
Since clear skies heavily outnumber EPB events in nocturnal observations, we faced a substantially imbalanced dataset. Left unchecked, a model would simply learn to predict a clear sky every time. We tackled this using a dual-strategy approach. During initial preprocessing, we oversample the minority class by applying a duplication multiplier to the images containing EPBs, being careful to preserve their structural integrity. Later, just before classifier training, we bring in the Synthetic Minority Over-sampling Technique (SMOTE) to interpolate artificial samples in the feature space. Together, these steps guarantee the algorithm trains on balanced data.

\subsection{Feature Extraction via 2DPCA}
Traditional Principal Component Analysis (PCA) requires flattening a 2D image into a long 1D vector, which inevitably tears apart the natural spatial correlations within the picture. Two-Dimensional Principal Component Analysis (2DPCA) avoids this problem entirely by working directly on the image matrices. Our tests showed that keeping just the first five principal components was entirely sufficient to capture the core visual signatures of EPBs. This cuts down the data size dramatically, lowering both the risk of overfitting and the overall processing time while preserving the spatial relationships that matter most for detection.

\subsection{Stacking Ensemble Classifier}
Rather than relying on a single algorithm to act as the core classification engine, we built a Stacking Ensemble architecture. This lets us borrow the unique strengths of several different learners at once. We train three distinct base classifiers side-by-side on the 2DPCA-extracted features:
\begin{itemize}
    \item \textbf{Support Vector Classifier (SVC):} Uses a Radial Basis Function (RBF) kernel to draw complex decision boundaries.
    \item \textbf{XGBoost Classifier:} Brought in for its strong non-linear feature extraction and gradient boosting capabilities.
    \item \textbf{Random Forest Classifier:} Reduces variance by building an ensemble of decision trees based on random subspaces.
\end{itemize}
The predictions from these three models feed into a Logistic Regression meta-classifier, which essentially learns how to weigh their output to make the final call (EPB present or absent). Because our preprocessing and feature extraction steps complemented each other so well, the architecture reached strong classification performance right out of the box, letting us skip exhaustive hyperparameter tuning entirely.

\subsection{Model Interpretability and Validation}
We wanted to be certain the model was actually "seeing" physical features, rather than just exploiting random artifacts like camera borders. To achieve this, we integrated SHapley Additive exPlanations (SHAP) to map out exactly which regions of an image drive the classifier's decisions. We also tested the model over continuous clear-sky sequences to ensure its prediction probabilities hold steady during long operational shifts.

\subsection{Evaluation Metrics}
Given the original imbalance in the data, relying on just one metric wouldn't tell the whole story. We evaluated the model using Accuracy, Sensitivity (Recall), Specificity, and the F1-Score. Sensitivity tells us how many actual EPBs we caught, ensuring we don't miss critical events. Specificity tracks how well we identified clear skies, keeping false positives in check. Finally, the F1-Score blends precision and sensitivity to give a clear, balanced picture of the model's overall robustness.

\chapter{Results}\label{chap:results}

The performance and robustness of the proposed Equatorial Plasma Bubble (EPB) detection architecture were evaluated through statistical validation, model interpretability analysis, and continuous temporal application. 

\subsection{Performance Metrics and Cross-Validation}
To make sure the model could generalize and wasn't just memorizing the training data, we ran a 5-fold cross-validation. The entire pipeline, from the 2DPCA feature extraction to the stacking ensemble and SMOTE, performed remarkably consistently across all data partitions. In fact, four out of the five folds hit perfect scores (1.0000) across Accuracy, Sensitivity, Specificity, and F1-Score. The third fold dipped slightly, logging an Accuracy of 0.9902, a Sensitivity of 0.9804, a Specificity of 1.0000, and an F1-Score of 0.9901. Averaging these out, the cross-validation yielded an Accuracy and F1-Score of $99.80\% \pm 0.44\%$. These numbers suggest the system can identify plasma structures with near-perfect reliability on this dataset.

\subsection{Model Interpretability}
We used the SHAP TreeExplainer to investigate the physical relevance of the features driving these decisions. Heatmap analyses showed that when an image actually contained an EPB, the model almost always focused its attention on the upper half of the optical field of view. This makes perfect physical sense, as plasma bubbles typically appear as vertical depletion structures rising into the upper quadrants. The lower half of the images rarely influenced the final decision weight. Interestingly, we also found a very low correlation between the absolute SHAP values and the 2DPCA reconstruction error. This implies that the decision tree ensemble and the mathematical image reconstruction are actually looking at different underlying image properties to confirm a detection.

\subsection{Temporal Stability Analysis}
Beyond static classification, we put the model through a continuous temporal analysis to see how it would hold up during a long operational shift. We fed it a complete, uninterrupted sequence of nocturnal airglow observations from the imager located at Bom Jesus da Lapa - Bahia. The sequence was recorded on July 22, 2025, running from 21:56 to 08:23 UT. Across those 230 consecutive images of a quiet ionosphere (meaning zero EPB events), the model showed strong immunity to false positives. The average EPB probability hovered at just $0.5\%$ for the entire night, and the absolute maximum noise peak only reached $17.7\%$, staying safely below the threshold that would trigger a classification alert. These findings confirm that the model is stable and reliable enough for real-time monitoring.

\chapter{Schedule}\label{chap:schedule}

To ensure the successful execution of this dissertation, we have organized the workload over eight quarters, as outlined in Table \ref{tab:schedule}. We will begin with a systematic literature review that will continuously run from the first quarter through to the eighth. The database construction, preprocessing, and model implementations will take place through the middle quarters, allowing us to transition into evaluation and writing phases smoothly.

\begin{table}[htb]
    \centering
    \caption{Dissertation Execution Schedule over 8 Quarters}
    \label{tab:schedule}
    \begin{tabular}{p{5.5cm} *{8}{c}}
        \toprule
        Activities & Q1 & Q2 & Q3 & Q4 & Q5 & Q6 & Q7 & Q8 \\
        \midrule
        Systematic Literature Review & X & X & X & X & X & X & X & X \\
        \addlinespace
        Labeled Database Construction &  & X & X & X &  &  &  &  \\
        \addlinespace
        Preprocessing and Data Augmentation &  &  & X & X &  &  &  &  \\
        \addlinespace
        Machine Learning Models Implementation &  &  & X & X & X & X &  &  \\
        \addlinespace
        Explainable AI Experiments &  &  &  & X & X & X &  &  \\
        \addlinespace
        Quantitative and Qualitative Evaluation &  &  &  &  & X & X &  &  \\
        \addlinespace
        \midrule
        Methodology and Results Chapters Drafting &  &  &  & X & X & X &  &  \\
        \addlinespace
        Critical Analysis and Results Discussion &  &  &  &  & X & X & X &  \\
        \addlinespace
        Qualification Defense Preparation &  &  &  &  &  & X &  &  \\
        \addlinespace
        Post Qualification Review and Adjustments &  &  &  &  &  &  & X &  \\
        \addlinespace
        Article Submission for Conference or Journal &  &  &  &  &  & X & X &  \\
        \addlinespace
        Final Dissertation Drafting and Formatting &  &  &  &  &  &  & X & X \\
        \addlinespace
        Final Dissertation Defense &  &  &  &  &  &  &  & X \\
        \bottomrule
    \end{tabular}
\end{table}

\chapter{Conclusion}\label{chap:conclusion}

The preliminary results presented in this qualification demonstrate that grounding image analysis in the physical characteristics of the observed phenomenon, rather than relying on brute-force feature learning, leads to a detection pipeline that is both accurate and transparent. The preprocessing stage effectively isolates the plasma depletion signatures from environmental contamination, and the compact feature representation produced by 2DPCA retains enough structural information for the ensemble classifier to distinguish between EPB events and clear-sky conditions with high confidence. The interpretability analysis further reinforced this outcome by confirming that the classifier's attention aligns with the expected spatial distribution of ionospheric depletions.

These findings establish a solid foundation for the next phases of this dissertation. Moving forward, a central priority will be assessing how well this architecture generalizes to observations from different geographical stations and across varying solar cycle conditions, where EPB morphology and occurrence rates differ considerably. We also plan to investigate the integration of complementary data sources, such as GNSS scintillation indices and radar backscatter measurements, to provide a more comprehensive view of ionospheric irregularities. Finally, exploring alternative and more recent machine learning architectures will allow us to benchmark the current approach and identify potential improvements for operational deployment.


% ================= POST-TEXTUAL =================
\postextual

% --- Referências ---
\bibliography{sections/biblio}

\end{document}
